% ============================================================
% APPENDIX — Production Artifacts and Workflow Documentation
% ============================================================

\appendix

% [R2-56: Appendix scope narrowed to OPM-specific content only]
\section{Plugin Architecture: The Open Paper Machine}
\label{app:plugin}

This appendix documents the \textit{Open Paper Machine specifically}---its plugin structure and orchestration dialogue. For general background on transformer architecture, RLHF, and Constitutional AI, see Section~\ref{sec:tech}.

The Open Paper Machine is a Claude Code plugin comprising 3{,}938 lines of prompt engineering, orchestration logic, and conversion utilities distributed across 18 files. The complete source is archived at \url{https://github.com/TobiasBlask/open-paper-machine}.

\subsection{System Overview}

The plugin consists of four architectural layers:

\begin{enumerate}
  \item \textbf{Plugin manifest} (\texttt{plugin.json}, 18 lines). Declares two MCP server dependencies: \texttt{academic-search} (Semantic Scholar, OpenAlex, CrossRef, arXiv APIs) and \texttt{paperbanana} (AI figure generation via Google Gemini). Both servers start automatically when Claude Code loads the plugin.

  \item \textbf{Autonomous agent} (\texttt{agents/paper-machine.md}, 570 lines). The core orchestration prompt that drives the six-phase pipeline. Operates on five principles: (1)~\textit{Do, don't ask}---make decisions and present results; (2)~\textit{Produce text, not plans}---every phase yields deliverable output; (3)~\textit{Checkpoint, don't block}---present work, then continue; (4)~\textit{Be explicit about decisions}---state what was chosen and why; (5)~\textit{Save everything to files}---every phase produces artifacts.

  \item \textbf{Skill engines} (7 modules, 1{,}933 lines total). Specialized execution modules activated on demand by the agent or user commands. Each encapsulates domain knowledge via fill-in-the-blank templates rather than abstract advice. Table~\ref{tab:skills} summarizes all engines.

  \item \textbf{Command interfaces} (8 commands, 685 lines total). User-facing entry points (\texttt{/write-paper}, \texttt{/search-papers}, \texttt{/export-latex}, \texttt{/verify-citations}, \texttt{/draft-section}, \texttt{/generate-figure}, \texttt{/generate-plot}, \texttt{/respond-reviewers}) that activate specific skill engines or the full pipeline.
\end{enumerate}

\begin{table}[H]
\centering
\small
\caption{Skill engines of the Open Paper Machine plugin.}
\label{tab:skills}
\begin{tabular}{@{}llp{7.5cm}@{}}
\toprule
\textbf{Engine} & \textbf{Lines} & \textbf{Capability} \\
\midrule
writing-engine & 339 & Sentence, paragraph, and section templates for IS/WI/BWL research. Provides concrete fill-in-the-blank formulas (e.g., 6-paragraph introduction blueprint, IMRAD section patterns). \\
literature-engine & 213 & API-first literature discovery across four academic databases. Search strategy generation, concept matrix construction, citation snowballing, narrative synthesis. \\
theory-engine & 169 & Theory selection (TOE, UTAUT, Dynamic Capabilities, etc.), research gap formulation, hypothesis derivation, contribution statement templates. \\
method-engine & 353 & Methodology decision trees for SLR, case study, Gioia, Mayring, grounded theory, PLS-SEM, DSR, and mixed methods. Complete method section templates. \\
figure-engine & 202 & AI figure generation via PaperBanana MCP (Gemini backend) with matplotlib/seaborn fallback. Transforms text descriptions into publication-quality diagrams. \\
latex-engine & 281 & Markdown-to-LaTeX conversion, \texttt{\textbackslash citep}/\texttt{\textbackslash citet} resolution, booktabs tables, arxiv-style compilation. Uses \texttt{md\_to\_latex.py} (732 lines). \\
verification-engine & 376 & Citation verification against source abstracts and full-text. Classifies each citation as VERIFIED, PLAUSIBLE, MISMATCH, UNVERIFIABLE, or NOT FOUND. \\
\bottomrule
\end{tabular}
\end{table}

\subsection{Artifact Inventory}

The pipeline produces the following intermediate artifacts, all version-controlled in the repository:

\begin{itemize}
  \item \texttt{literature\_base.csv} --- Retrieved literature database (Phase 1)
  \item \texttt{references.bib} --- BibTeX entries for all cited papers (Phase 1)
  \item \texttt{framing.md} --- Theory selection, gap formulation, research questions (Phase 2)
  \item \texttt{concept\_matrix.md} --- Webster \& Watson concept matrix (Phase 3)
  \item \texttt{paper\_structure.md} --- Section structure with word budget (Phase 3)
  \item \texttt{draft.md} --- Complete paper draft as Markdown (Phase 4)
  \item \texttt{figures/} --- PaperBanana-generated diagrams, 300 DPI PNG (Phase 5)
  \item \texttt{latex/paper.tex} --- arxiv-style LaTeX source (Phase 5)
  \item \texttt{latex/paper.pdf} --- Compiled PDF manuscript (Phase 5)
  \item \texttt{status\_report.md} --- Self-assessed completion status (Phase 6)
  \item \texttt{verification\_report.md} --- Citation verification results (optional Phase 7)
\end{itemize}

% [R3: Appendix B (Prototypical Dialogue) deleted per Burkhardt — orchestrator role appeared too passive in the example]

% ============================================================
\section{Orchestrator--Agent Interaction Timeline}
\label{app:timeline}

This appendix documents the complete sequence of substantive orchestrator--agent interactions during the production of this paper. Each entry records the pipeline phase, the \textit{character} of the interaction (approval, redirect, expansion, correction, or evaluation), and what it reveals about the orchestrator role. Routine interactions (``proceed,'' ``looks good'') are omitted; only interactions that shaped the paper's content or direction are included.

\vspace{0.5em}

\begin{table}[H]
\centering
\small
\caption{Orchestrator--agent interaction timeline. Character types: \textbf{D}~= Direction (new task or redirect), \textbf{E}~= Expansion (add scope), \textbf{C}~= Correction (fix error), \textbf{V}~= Evaluation (quality judgment), \textbf{A}~= Approval (gate passed).}
\label{tab:timeline}
\begin{tabularx}{\textwidth}{@{} c l c X @{}}
\toprule
\textbf{\#} & \textbf{Phase} & \textbf{Char.} & \textbf{Interaction} \\
\midrule
1 & Phase 1 & D & \textit{Orchestrator} initiates pipeline with title and research topic. Agent autonomously executes 6 search queries across 4 databases, returns 171 candidates. \\[0.3em]
2 & Phase 1 & E & \textit{Orchestrator} judges sociotechnical theory coverage as thin. Commissions 3 additional targeted searches (``sociotechnical envelopment AI,'' ``task-based automation framework,'' ``distributed cognition AI systems''), adding 14 papers. \\[0.3em]
3 & Phase 1 & A & Corpus approved: 62 papers retained for concept matrix. \\[0.3em]
4 & Phase 2 & D & \textbf{Key redirect.} Agent frames paper as pure systems description. \textit{Orchestrator rejects} and redirects toward creator-to-orchestrator thesis---the substantive intellectual decision that shaped the entire argument. \\[0.3em]
5 & Phase 2 & V & \textit{Orchestrator} evaluates theory selection (distributed cognition, sociotechnical systems, Autor). Approves combination; adds Scarbrough (2024) as missing empirical complement. \\[0.3em]
6 & Phase 3 & V & \textit{Orchestrator} reviews section structure. Rebalances word budget: reduces technology description (\S4), expands discussion (\S6). \\[0.3em]
7 & Phase 3 & A & Structure approved: 8 sections, $\sim$10{,}000 words. \\[0.3em]
8 & Phase 4 & C & \textit{Orchestrator} spot-checks citations against retrieved abstracts. Identifies 4 hallucinated claims (correct paper, incorrect attributed finding). Agent corrects. \\[0.3em]
9 & Phase 4 & V & \textit{Orchestrator} evaluates prose quality. Flags ``intellectual risk'' deficit: too hedged, too balanced. Instructs agent to sharpen provocative arguments (esp.\ \S6.2 desk research). \\[0.3em]
10 & Phase 5 & C & \textbf{Figure correction.} Co-author identifies conceptual error: ``accountability'' arrow from AI System to Published Output contradicts core argument. Agent confirms error, regenerates figure. \\[0.3em]
11 & Phase 5 & A & \LaTeX{} compilation approved. All figures, tables, references resolve. \\[0.3em]
12 & Phase 6 & V & Agent generates self-review (simulated peer review). \textit{Orchestrator} evaluates 12 critique points: accepts 8, rejects 4 as unfounded. Accepted critiques drive final revision pass. \\[0.3em]
13 & Co-author & D & Co-author review introduces new requirements: position paper framing, shorter \S4, ``desk research $\rightarrow$ desk reject'' argument. Agent implements; orchestrator verifies. \\[0.3em]
14 & R2 Revision & E & 59 review comments processed. Agent drafts responses and edits; orchestrator evaluates each for argumentative coherence and reviewer intent. \\
\bottomrule
\end{tabularx}
\end{table}

\subsection{Patterns in the Interaction}

Three patterns emerge from the timeline:

\begin{enumerate}
    \item \textbf{The orchestrator's most consequential interventions are redirections, not corrections.} Interaction \#4 (rejecting the systems framing in favor of the creator-to-orchestrator thesis) shaped the entire paper more than any error correction. This is the ``step up'' strategy \citep{davenport2016just}: the orchestrator operates at the level of research direction, not prose production.

    \item \textbf{Corrections cluster at the semantic level, not the syntactic level.} The agent produces grammatically correct, convention-following prose. Errors occur at the level of \textit{meaning}: hallucinated claims (\#8), conceptual contradictions in figures (\#10), insufficient intellectual risk (\#9). These are precisely the errors that require domain expertise to detect.

    \item \textbf{The interaction character shifts across phases.} Early phases are dominated by \textit{direction} and \textit{expansion} (what to research, how to frame it). Middle phases shift to \textit{correction} and \textit{evaluation} (are the claims right? is the prose sharp enough?). Late phases return to \textit{direction} as co-author feedback introduces new requirements. This pattern---\textit{direct} $\rightarrow$ \textit{evaluate} $\rightarrow$ \textit{redirect}---mirrors the orchestrator cycle depicted in Figure~\ref{fig:creator_to_orchestrator}.
\end{enumerate}

\subsection{What the Timeline Reveals}

The timeline operationalizes the ``orchestration is substantive intellectual work'' claim from Section~3.2. Across 14 substantive interactions, the orchestrator:
\begin{itemize}[nosep]
    \item Made 3 directional decisions that shaped the paper's argument (\#1, \#4, \#13)
    \item Expanded scope based on domain knowledge 2 times (\#2, \#14)
    \item Corrected semantic errors the agent could not self-detect 2 times (\#8, \#10)
    \item Exercised evaluative judgment on quality, balance, and sharpness 4 times (\#5, \#6, \#9, \#12)
\end{itemize}

\noindent None of these interactions involved text production. All involved judgment. This is the creator-to-orchestrator transition made observable.

% End of appendix content
